% ============================================================================
% EXEMPLO - GRADIENTES NEWTON PAIVA
% Template demonstrando diferentes estilos de gradiente
% ============================================================================
\documentclass[aspectratio=169]{beamer}

% Usa o sistema de layouts Newton Paiva
\usepackage{../styles/newton-paiva-layouts}

% Pacotes úteis
\usepackage[utf8]{inputenc}
\usepackage[portuguese]{babel}

% ============================================================================
% CONFIGURAÇÕES DA APRESENTAÇÃO
% ============================================================================

% Título e subtítulo
\title{Demonstração de Gradientes}
\subtitle{Planos de fundo personalizados Newton Paiva}

% Informações da Newton Paiva
\disciplina{Design de Apresentações}
\instituicao{Centro Universitário Newton Paiva}
\professor{Professor Exemplo}

% Data
\date{\today}

% Seção atual
\secaoatual{Gradientes}

% ============================================================================
% INÍCIO DO DOCUMENTO
% ============================================================================
\begin{document}

% ============================================================================
% TÍTULO COM GRADIENTE HORIZONTAL (PADRÃO)
% ============================================================================
\titlepagepadrao
\begin{frame}[plain]
    \maketitle
\end{frame}

% ============================================================================
% SLIDES COM DIFERENTES GRADIENTES
% ============================================================================

\layoutpadrao

\begin{frame}{Gradiente Horizontal (Padrão)}
    \begin{block}{Características}
        \begin{itemize}
            \item Gradiente da esquerda para direita
            \item Roxo/magenta → Laranja/vermelho
            \item Baseado na imagem fornecida
            \item Ideal para a maioria dos slides
        \end{itemize}
    \end{block}
    
    \begin{exampleblock}{Comando}
        \texttt{\textbackslash fundogradiente} (já ativo por padrão)
    \end{exampleblock}
\end{frame}

% Mudando para gradiente vertical
\fundogradientevertical

\begin{frame}{Gradiente Vertical}
    \begin{center}
        \vfill
        \Huge{\textcolor{white}{\textbf{Gradiente Vertical}}}
        \vspace{1cm}
        
        \Large{\textcolor{white}{De cima para baixo}}
        \vspace{1cm}
        
        \begin{itemize}
            \item Roxo/magenta no topo
            \item Laranja/vermelho na base
            \item Efeito dramático
            \item Ideal para slides de seção
        \end{itemize}
        
        \vspace{1cm}
        \texttt{\textcolor{yellow}{\textbackslash fundogradientevertical}}
        \vfill
    \end{center}
\end{frame}

% Mudando para gradiente radial
\fundogradienteradial

\begin{frame}{Gradiente Radial}
    \begin{center}
        \vfill
        \Huge{\textcolor{white}{\textbf{Gradiente Radial}}}
        \vspace{1cm}
        
        \Large{\textcolor{white}{Do centro para as bordas}}
        \vspace{1cm}
        
        \begin{itemize}
            \item Centro com cor intermediária
            \item Bordas com roxo/magenta
            \item Efeito de profundidade
            \item Destaca o conteúdo central
        \end{itemize}
        
        \vspace{1cm}
        \texttt{\textcolor{yellow}{\textbackslash fundogradienteradial}}
        \vfill
    \end{center}
\end{frame}

% Voltando para fundo sólido
\fundosolido

\begin{frame}{Fundo Sólido Original}
    \begin{alertblock}{Comparação}
        Este é o fundo sólido original da Newton Paiva para comparação.
    \end{alertblock}
    
    \begin{block}{Comandos Disponíveis}
        \begin{itemize}
            \item \texttt{\textbackslash fundogradiente} - Gradiente horizontal (padrão)
            \item \texttt{\textbackslash fundogradientevertical} - Gradiente vertical
            \item \texttt{\textbackslash fundogradienteradial} - Gradiente radial
            \item \texttt{\textbackslash fundosolido} - Fundo sólido original
        \end{itemize}
    \end{block}
    
    \begin{exampleblock}{Uso}
        Coloque o comando antes do slide que quer alterar o fundo.
    \end{exampleblock}
\end{frame}

% Voltando para o gradiente horizontal
\fundogradiente

\begin{frame}{Dicas de Uso}
    \begin{columns}
        \begin{column}{0.5\textwidth}
            \begin{block}{Gradiente Horizontal}
                \begin{itemize}
                    \item Slides normais
                    \item Conteúdo geral
                    \item Leitura fácil
                \end{itemize}
            \end{block}
            
            \begin{block}{Gradiente Vertical}
                \begin{itemize}
                    \item Slides de título de seção
                    \item Apresentações dramáticas
                    \item Destaque visual
                \end{itemize}
            \end{block}
        \end{column}
        \begin{column}{0.5\textwidth}
            \begin{exampleblock}{Gradiente Radial}
                \begin{itemize}
                    \item Slides de foco central
                    \item Conclusões importantes
                    \item Chamadas de atenção
                \end{itemize}
            \end{exampleblock}
            
            \begin{alertblock}{Fundo Sólido}
                \begin{itemize}
                    \item Compatibilidade máxima
                    \item Impressão em preto/branco
                    \item Ambiente conservador
                \end{itemize}
            \end{alertblock}
        \end{column}
    \end{columns}
\end{frame}

\begin{frame}
    \begin{center}
        \vfill
        \Huge{\textcolor{white}{\textbf{Obrigado!}}}
        \vspace{1cm}
        
        \newtonlogo[3cm]
        
        \vspace{1cm}
        \Large{\textcolor{white}{Gradientes Newton Paiva}}
        
        \vspace{1.5cm}
        \textcolor{newtanorange}{\textbf{Sistema de Apresentações}}
        \\[0.3cm]
        \textcolor{white}{Centro Universitário Newton Paiva}
        \vfill
    \end{center}
\end{frame}

\end{document}
